\section{Worked Inconsistency and Repair Cases}\label{app:workedexamples}

The following four recorded benchmark cases connect the original artifacts, verification evidence,
repair and subsequent checks. They connect individual mechanisms to exact variable domains,
objective directions and constraint coefficients.

\subsection{A ReLoop pass with inconsistent feasible regions}
\label{app:industryexample}

\input{generated_consistency/case.tex}

IndustryOR problem 71 at 14B returns the reference value 623. The mathematical model and program
each contain four variables and five constraints. ReLoop accepts the program's execution, and the
reported solution passes feasibility, value binding and optimality checks under both models.
The models nevertheless permit different decisions. Figure~\ref{fig:relooppass} shows the original
judgment, the two witnesses and the two repair directions.

\textbf{Complete compared models.} Let \(b,p,r,t\) denote boat, plane, train and truck quantities.
All four are nonnegative integers. Both artifacts maximize
\begin{equation}
5t+10p+8b+7r
\label{eq:industryobjective}
\end{equation}
and contain the four common constraints
\begin{equation}
\begin{aligned}
12t+20p+15b+10r&\leq890,\\
6t+3p+5b+4r&\leq500,\\
t-r&\leq0,\\
b-r&\leq0.
\end{aligned}
\label{eq:industrycommon}
\end{equation}
The remaining mathematical row is \(p-b\geq0\), while the program adds \(b-p\geq0\).
Their objective coefficients, row counts and variable correspondence agree. The different
inequality restricts the feasible region in both directions.

\textbf{Witness checks.} The decision \((b,p,r,t)=(0,1,0,0)\) satisfies every mathematical
constraint. Its wood and steel use are 20 and three, and its executable boat row has value negative
one. Conversely, \((1,0,1,0)\) satisfies the program's complete model, uses 25 units of wood and nine
of steel, and violates the written boat row by one. These decisions are witnesses to the two
set differences, independently of the optimum shared by the models.

\begin{table}[htbp]
\caption{Exact witness evaluations for IndustryOR 71 at 14B.}
\label{tab:industrywitnesses}
\centering\small\setlength{\tabcolsep}{4pt}
\begin{tabular*}{\textwidth}{@{\extracolsep{\fill}}lrr@{}}
\toprule
Quantity or check & Mathematical witness & Executable witness\\
\midrule
Decision \((b,p,r,t)\) & \((0,1,0,0)\) & \((1,0,1,0)\)\\
Wood use, limit 890 & 20 & 25\\
Steel use, limit 500 & 3 & 9\\
Mathematical row \(p-b\), lower bound 0 & 1 & \(-1\)\\
Executable row \(b-p\), lower bound 0 & \(-1\) & 1\\
Objective value & 10 & 15\\
All mathematical constraints satisfied & Yes & No\\
All executable constraints satisfied & No & Yes\\
\bottomrule
\end{tabular*}
\end{table}

\textbf{Repair from the mathematical model.} Keeping the mathematical model as the selected
target replaces the executable inequality with \(p-b\geq0\). The compiler also reproduces the
four common rows, integer declarations and objective. Capturing the new program yields the same
complete model. Both difference queries are unsatisfiable, and the objective expressions are equal.
This direction repairs the inconsistent implementation from the available mathematics.

\textbf{Repair from the program.} Keeping the captured program as the selected target exports
\(b-p\geq0\) into the mathematical artifact and retains every other captured field. The regenerated
program and exported model again have unsatisfiable difference queries and equal objectives.
This direction completes or corrects mathematics from an executable model. The source requirements
determine which selected interpretation should be retained.

\begin{evidencecase}{IndustryOR 71: the changed code and the certified result}
\begin{minipage}[t]{0.475\linewidth}\vspace{0pt}\small
\textcolor{verdictblue}{\textbf{Original program row.}}\par\smallskip
\texttt{m.addConstr(boat - plane >= 0)}\par\medskip
\textcolor{verdictgreen}{\textbf{Mathematical target repair.}}\par\smallskip
\texttt{m.addConstr(plane - boat >= 0)}\par\medskip
The variable declarations, objective and four other rows are retained. The program prints its
computed objective and variable assignments after optimization.
\end{minipage}\hfill
\begin{minipage}[t]{0.475\linewidth}\vspace{0pt}\small
\textcolor{verdictgreen}{\textbf{Certification after construction.}}\par\smallskip
Both region difference queries are unsatisfiable. The objective coefficients and direction match.
A decision with value 623 exists, and the query for a value greater than 623 is unsatisfiable.\par\medskip
\textcolor{verdictblue}{\textbf{Effect of the repair.}}\par\smallskip
The model--code relation changes from refuted to certified even though the answer remains correct.
\end{minipage}
\end{evidencecase}

\textbf{Executed reconstruction.} A fresh replay executes the original program and both
reconstructed Gurobi programs. It reads the instantiated variables, bounds, rows and objective from
each program and compares them with the selected target using exact Z3 queries. Both repaired
programs return \((b,p,r,t)=(0,0,89,0)\) with value 623. The returned point satisfies the complete
model and binds to that value, and no better feasible decision exists. This replay checks the
actual constructed programs in addition to the displayed row calculations.

\subsection{Missing source requirements in two generated artifacts}
\label{app:mamoexample}

MAMO easy problem 20 at 1.5B illustrates two distinct repairs. The first restores agreement between
the generated mathematical model and code. The second restores source requirements that both
artifacts omitted and changes the optimal result.

\textbf{Source and generated models.} The problem chooses integer quantities of four tour
packages, with costs 100, 200, 300 and 400. It gives two capacity limits and two minimum quality
requirements. The generated model uses nonnegative integers \(x_1,x_2,x_3,x_4\), minimizes
\begin{equation}
f(x)=100x_1+200x_2+300x_3+400x_4,
\label{eq:mamocost}
\end{equation}
and contains only \(x_1+x_2\leq50\). The program also includes \(x_3+x_4\leq80\).
Both artifacts omit the quality requirements. Their reported decision is zero in every coordinate,
with objective zero. The recorded report has no explicit solver status.

\begin{evidencecase}{MAMO easy 20: disagreement and missing requirements}
\begin{minipage}[t]{0.475\linewidth}\vspace{0pt}\small
\textcolor{verdictblue}{\textbf{Mathematical artifact.}}\par
\(x_i\in\mathbb Z_{\geq0}\), minimize \(f(x)\),\par
\(x_1+x_2\leq50\).\par\medskip
\textcolor{verdictorange}{\textbf{Additional code row.}}\par
\(x_3+x_4\leq80\).\par\medskip
\textcolor{verdictorange}{\textbf{Exact counterexample.}}\par
\((0,0,81,0)\) satisfies the written model and violates the additional capacity row by one.
\end{minipage}\hfill
\begin{minipage}[t]{0.475\linewidth}\vspace{0pt}\small
\textcolor{verdictblue}{\textbf{Omitted source requirements.}}\par
\(0.6x_1+0.8x_2\geq20\),\par
\(0.7x_3+0.9x_4\geq30\).\par\medskip
\textcolor{verdictorange}{\textbf{Old report.}}\par
\((0,0,0,0)\) has zero cost and is optimal under both incomplete artifacts, while violating both
source quality requirements.
\end{minipage}
\end{evidencecase}

\textbf{Artifact repair.} Adding the second capacity row to the written model makes its region
identical to the program's region. The objective and domains were equal before this change and
remain equal. The zero decision remains optimal. This repair resolves the artifact discrepancy
without changing the interpretation shared by both generated models.

\textbf{Source repair.} The reviewed source formulation contains both capacities and both
quality requirements. Clearing the decimal denominators gives the complete repaired model
\begin{equation}
\begin{aligned}
\min\quad &100x_1+200x_2+300x_3+400x_4\\
\text{subject to}\quad
 &x_1+x_2\leq50, &x_3+x_4&\leq80,\\
 &3x_1+4x_2\geq100, &7x_3+9x_4&\geq300,\\
 &x_1,x_2,x_3,x_4\in\mathbb Z_{\geq0}.&&
\end{aligned}
\label{eq:mamosourcerepair}
\end{equation}
The integer coefficients represent the source decimals exactly. The old zero report fails both
quality checks. The repaired decision \((34,0,43,0)\) has value 16,300 and matches the released
answer. A fresh exact query for a feasible decision with lower cost is unsatisfiable.

\begin{table}[htbp]
\caption{Feasibility and value checks for the repaired MAMO easy 20 decision.}
\label{tab:mamorepairchecks}
\centering\small\setlength{\tabcolsep}{4pt}
\begin{tabular*}{\textwidth}{@{\extracolsep{\fill}}lrr@{}}
\toprule
Constraint or objective & Evaluated value & Required relation\\
\midrule
First capacity, \(x_1+x_2\) & 34 & At most 50\\
Second capacity, \(x_3+x_4\) & 43 & At most 80\\
First quality, \(3x_1+4x_2\) & 102 & At least 100\\
Second quality, \(7x_3+9x_4\) & 301 & At least 300\\
Objective, \(f(x)\) & 16,300 & Equals reported value\\
Feasible decision with smaller objective & None & Query unsatisfiable\\
\bottomrule
\end{tabular*}
\end{table}

\textbf{Result of each intervention.} Model--code verification first restores the missing capacity
row, changing the mathematical/executable row counts from one/two to two/two and the model--code
relation from refuted to certified; the objective remains zero and the answer remains incorrect.
Source review then restores both quality rows in both artifacts, giving four/four rows, objective
16,300 and a correct answer. Code construction propagates the completed mathematical model into the
implementation; the returned decision passes the feasibility and value checks above.

\subsection{A domain mismatch and redundant constraints}
\label{app:complexexample}

ComplexOR problem 3 at 1.5B shows why row agreement alone is insufficient. The written model uses
nonnegative integers \(x,y\), maximizes \(x+y\), and repeats the same row three times:
\begin{equation}
c_1:x+y\leq2,\qquad c_2:x+y\leq2,\qquad c_3:x+y\leq2.
\label{eq:complexrows}
\end{equation}
The code uses the same objective and rows, but declares \(x\) binary. The decision \((2,0)\)
satisfies all rows of both artifacts and violates only the executable domain. The original report
\((0,2)\), with value two, is feasible, correctly valued and optimal under both models.

\begin{evidencecase}{ComplexOR 3: domain repair and a smaller equivalent constraint set}
\begin{minipage}[t]{0.475\linewidth}\vspace{0pt}\small
\textcolor{verdictorange}{\textbf{Domain difference.}}\par
Mathematics: \(x\in\mathbb Z_{\geq0}\).\par
Code: \(x\in\{0,1\}\).\par\medskip
\textcolor{verdictorange}{\textbf{Separating decision.}}\par
At \((x,y)=(2,0)\), each row has left hand side two, while the binary domain rejects \(x=2\).\par\medskip
\textcolor{verdictgreen}{\textbf{Domain repair.}}\par
Declare the code variable as a nonnegative integer and check both inclusions again.
\end{minipage}\hfill
\begin{minipage}[t]{0.475\linewidth}\vspace{0pt}\small
\textcolor{verdictblue}{\textbf{Redundancy evidence.}}\par
For each duplicate, the other two rows entail it. Delete \(c_3\), then test and delete \(c_2\)
against the remaining model.\par\medskip
\textcolor{verdictgreen}{\textbf{Final model.}}\par
One row \(c_1:x+y\leq2\) remains. The integer domains and objective are unchanged. The maximum
value remains two.
\end{minipage}
\end{evidencecase}

\textbf{Checked stages.} The original three-row mathematical model is the target, whereas the
three-row executable model is refuted because it makes \(x\) binary. Restoring \(x\) as a
nonnegative integer certifies the pair. Sequentially removing one and then two duplicate rows
retains certification, leaving one essential row with the same domains and objective.

\begin{figure}[htbp]
\begin{minipage}[t]{0.52\textwidth}\vspace{0pt}
\centering
\includegraphics[width=\linewidth]{generated_consistency/case_domain.pdf}
\caption{ComplexOR 3: the original binary domain excludes the circled integer decision; repair restores it.}
\label{fig:domaincase}
\end{minipage}\hfill
\begin{minipage}[t]{0.44\textwidth}\vspace{0pt}\small
\textbf{Sequential deletion.} To remove \(c_3\), the checker proves that \(c_1,c_2\) and
\(x+y>2\) cannot hold together. After deleting \(c_3\), the same query with \(c_1\) proves
\(c_2\) redundant. The remaining row is essential because \((3,0)\) violates the original bound.
Thus each deletion preserves an equivalent constraint set.

\smallskip
\textbf{Source interpretation.} The source concerns three participants, three cars and an interest
matrix, whereas both artifacts use two variables. Domain and redundancy repair certifies the
generated pair, but restoring the source requires indexed assignment variables and constraints.
The released answer is three; the repaired generated model still has optimum two.
\end{minipage}
\end{figure}

\subsection{An incorrect Opt-Verifier judgment on a certified solution}
\label{app:bworexample}

\begin{evidencecase}{BWOR 18: a correct result receives incorrect revision advice}
\begin{minipage}[t]{0.475\linewidth}\vspace{0pt}\small
\textcolor{verdictgreen}{\textbf{Returned assignment.}}\par
The companies choose stores 3, 2, 1, 4 and 5. The costs are 7, 9, 6, 6 and 6, giving 34 in total.\par\medskip
\textcolor{verdictgreen}{\textbf{Formal checks.}}\par
The model pair is consistent. The decision is feasible, its value is bound correctly, and no
better feasible decision exists in the compared models.
\end{minipage}\hfill
\begin{minipage}[t]{0.475\linewidth}\vspace{0pt}\small
\textcolor{verdictorange}{\textbf{Verifier comment.}}\par
The proposed alternative uses stores 1, 2, 3, 4 and 5. Its listed costs are 4, 9, 12, 6 and 6.
The comment labels this total as 34.\par\medskip
\textcolor{verdictblue}{\textbf{Arithmetic check.}}\par
The listed alternative costs total 37. The alternative supplies no improvement over the certified
value 34.
\end{minipage}
\end{evidencecase}

BWOR problem 18 at 14B assigns five construction companies to five stores. Its recorded model
contains 25 variables and ten constraints. Both region inclusions, objective coefficients and
preferences agree. The reported objective 34 matches the reference value, and the recorded exact
checks certify both region inclusions, equal objective coefficients, equal decision preferences,
feasibility under both models and objective binding. The query for a better feasible decision is
unsatisfiable.

The Opt-Verifier reproduction accepts the structural comparison but returns a negative solution
judgment. Its proposed alternative has store indices \((1,2,3,4,5)\) and costs
\((4,9,12,6,6)\), which sum to 37 rather than the claimed 34. The returned assignment
\((3,2,1,4,5)\) instead costs (7+9+6+6+6=34). The error is in the feedback arithmetic, not
model extraction or variable alignment.

\textbf{Exact decision path and consequence.} The checker proves both region inclusions and equal
objective coefficients, substitutes the returned assignment into every row, binds the reported
value to the evaluated objective and proves that no feasible assignment has value below 34. The
archived comparison judgment uses five language model calls and 44.2 seconds, yet its proposed
counterexample is strictly worse. The exact certificate therefore retains the correct program and
value instead of triggering an unnecessary revision.
